% ============================================================================
% 2.1 系统调用：最底层的 I/O
% ============================================================================
\subsection{为什么需要系统调用}

在用户态运行的程序，不能直接操作硬件。当我们想在屏幕上输出、从键盘读入，
或者结束程序，都必须请求内核代劳。这种请求就叫做\textbf{系统调用（system call）}。

在 ARM 64 位架构下，系统调用通过 \texttt{svc \#0x80} 指令触发。
操作系统收到这个中断后，根据 \texttt{x16} 寄存器里的编号，执行不同的操作。

\subsection{exit.s：三个系统调用的封装}

整个程序的所有 I/O 都由下面这几十个字节的代码完成：

\begin{lstlisting}[style=asm,caption=ISA/src/exit.s —— sys\_write / sys\_read / exit]
.global exit
.global sys_write
.global sys_read
.align 4

exit:
    movz x16, #0x1
    movk x16, #0x200, lsl #16
    svc #0x80
    ret

sys_write:
    movz x16, #0x4
    movk x16, #0x200, lsl #16
    svc #0x80
    ret

sys_read:
    movz x16, #0x3
    movk x16, #0x200, lsl #16
    svc #0x80
    ret
\end{lstlisting}

\subsection{代码解读}

\paragraph{编号拼接：movz + movk}
每个系统调用需要把编号放入 \texttt{x16}。ARM 上不能一步写入超过 16 位的立即数，
所以使用两步：

\begin{itemize}
    \item \texttt{movz x16, \#0x1} — 把 \texttt{x16} 的低 16 位置为 \texttt{0x1}，其余清零
    \item \texttt{movk x16, \#0x200, lsl \#16} — 保持其他位不变，把 16 位以上的部分写入 \texttt{0x200}
\end{itemize}

合起来就是 \texttt{0x2000001}（exit）、\texttt{0x2000003}（read）、
\texttt{0x2000004}（write），这是 macOS XNU 内核的系统调用编号约定。

\paragraph{参数传递}
三个函数都没有显式的参数处理，因为它们遵循 AArch64 的调用约定：
调用者把参数放在 \texttt{x0}、\texttt{x1}、\texttt{x2} 里，
\texttt{sys\_write / sys\_read / exit} 自己不需要再做什么——
只要把编号放到 \texttt{x16}，然后 \texttt{svc \#0x80} 就够了。
内核会自动从 \texttt{x0-x2} 取走参数。

举个例子——调用 \texttt{sys\_write} 向标准输出写一个字符串：

\begin{lstlisting}[style=asm,caption=sys\_write 的典型调用]
    mov  x0, #1           // 文件描述符：1 = 标准输出
    adrp x1, msg@PAGE     // 字符串所在页的基地址
    add  x1, x1, msg@PAGEOFF // 加页内偏移，得到真正的地址
    mov  x2, #msg_len     // 字符串长度
    bl   sys_write        // 调用 sys_write（bl = branch with link）
\end{lstlisting}

\paragraph{svc \#0x80：真正的内核切换}
这条指令的作用是：产生一个异常，CPU 从用户态切换到内核态，
跳转到内核中预先注册的异常向量表。内核根据 \texttt{x16} 的值执行对应处理函数，
完成后再返回用户态，继续执行下一条 \texttt{ret}。
如果没有这条指令，我们写的程序根本无法和外部世界交互。
